% ============================================================================
% INTRODUCCIÓN
% ============================================================================
% La introducción es la presentación clara, breve y precisa del contenido
% del trabajo. NO debe incluir resultados ni conclusiones.
%
% Según la norma UTEM, debe comprender:
%   a) Las razones que motivaron la elección del tema
%   b) La justificación: fundamentos y relevancia del trabajo
%   c) Formulación clara del problema, objetivos generales y naturaleza del trabajo
%   d) Un resumen de cómo se lograron los objetivos planteados
%   e) Una orientación al lector de cómo se ha organizado el texto
% ============================================================================

\chapter*{INTRODUCCIÓN}
\addcontentsline{toc}{chapter}{INTRODUCCIÓN}  % Agrega la Introducción a la tabla de contenido

% --- Contexto y motivación ---
% Explica las razones que motivaron la elección del tema y el contexto general.

[Escribe aquí el contexto y la motivación de tu trabajo...]

% --- Justificación ---
% Explica los fundamentos que sustentan el tema y la relevancia del trabajo.

% --- Problema de investigación ---
% Formula de manera clara el problema que investigaste, presentando los
% objetivos generales y la naturaleza del trabajo.

% --- Objetivos ---
% Describe brevemente los objetivos de tu trabajo.

% Objetivo general:
% - ...

% Objetivos específicos:
% - ...
% - ...
% - ...

% --- Estructura del documento ---
% Orienta al lector sobre cómo se ha organizado el texto.
% Ejemplo:
% "El presente trabajo se organiza de la siguiente manera: en el Capítulo 1 se
% presenta el marco teórico; en el Capítulo 2 se describe la metodología
% utilizada; en el Capítulo 3 se desarrolla la propuesta; en el Capítulo 4
% se presentan los resultados obtenidos; y finalmente, se exponen las
% conclusiones del trabajo."
